\begin{theorem}
Let $n\ge1$ be an integer and let $1<p<\infty$.  Write
$\mathbb{B}=\{\xi:|\xi|<1\}$ and $d\sigma$ for plane Lebesgue measure.  Then
there is a constant $C=C(n,p)>0$ with the following property.

Let $H\colon\mathbb{C}\to\mathbb{C}$ satisfy
\begin{enumerate}
\item[(H1)] $H\in W^{1,2}_{\mathrm{loc}}(\mathbb{B})$ in the sense of
  \noderef{MemWLocOn}, $H$ continuous on $\mathbb{B}$, and $H$ differentiable
  at almost every point of $\mathbb{B}$;
\item[(H2)] $H$ solves equation (18) on $\mathbb{B}$ in the sense of
  \noderef{SolvesEq18};
\item[(H3)] $|H_\xi|^2+|H_{\bar\xi}|^2$ integrable on $\mathbb{B}$, where
  $H_\xi$, $H_{\bar\xi}$ are the Wirtinger derivatives of \noderef{DZ} and
  \noderef{DZbar}.
\end{enumerate}
Then there are complex numbers $A_k$, $k\in\mathbb{Z}$, such that, with
$\lambda_k=\lambda_{|k|}(n,p)$ the exponent of \noderef{LambdaCoef},
$\epsilon_k=\epsilon_k(n,p)$ the coefficient of \noderef{EpsCoef}, $a_k$ the
Fourier coefficient of \noderef{FourierCoefOfH} and $T_k$ the $k$-th summand of
(19) in the sense of \noderef{Thm2Term}:
\begin{enumerate}
\item[(39)] $A_{-k}=\epsilon_k\overline{A_k}$ for every integer $k\ge n$;
\item[(b)] $a_k(r)=A_k\,r^{\lambda_k}$ for every $k\in\mathbb{Z}$ with
  $|k|\ge n$ and every $r\in(0,1)$;
\item[(22)] the series $\sum_{m\ge0}(n+1+m)|A_{n+1+m}|^2$ converges, and
\[
  \sum_{m=0}^{\infty}(n+1+m)\,|A_{n+1+m}|^2
  \le C\iint_{\mathbb{B}}\bigl(|H_\xi|^2+|H_{\bar\xi}|^2\bigr)d\sigma ,
\]
that is $\sum_{k=n+1}^{\infty}k|A_k|^2<\infty$ and
$\sum_{k=n+1}^{\infty}k|A_k|^2\le C\iint_{\mathbb{B}}(|H_\xi|^2+|H_{\bar\xi}|^2)d\sigma$;
\item[(19)] the series (19) converges to $H$ in $W^{1,2}(\mathbb{B})$:
\[
  \iint_{\mathbb{B}}\Bigl|H-\sum_{k=n}^{n+N}T_k\Bigr|^2d\sigma
  +\iint_{\mathbb{B}}\Bigl|H_\xi-\Bigl(\sum_{k=n}^{n+N}T_k\Bigr)_\xi\Bigr|^2d\sigma
  +\iint_{\mathbb{B}}\Bigl|H_{\bar\xi}
      -\Bigl(\sum_{k=n}^{n+N}T_k\Bigr)_{\bar\xi}\Bigr|^2d\sigma
  \longrightarrow0
\]
as $N\to\infty$.
\end{enumerate}

This is the forward half of Theorem 2 of the paper (paper.tex:187--205).  Its
hypothesis surface is now the printed one, $H\in W^{1,2}(\mathbb{B})$ solving
(18), together with two clauses that cost nothing: continuity of $H$ on
$\mathbb{B}$, which the inverse of a quasiconformal homeomorphism has because
it is a homeomorphism, and differentiability almost everywhere, which item (C2)
of the paper's note supplies for that same inverse and which is needed because
(18) and the polar derivatives are written with the Fr\'echet derivative.  In
particular the $C^\infty$ half of the normalisation licensed at paper.tex:183,
``$H\in C^\infty(\mathbb{B}-\{0\})$'', is no longer assumed: it is a
\emph{consequence} of the expansion: it is produced downstream by term-by-term
differentiation of (19), exactly as at paper.tex:387--388, and so cannot be
assumed here.  In accordance with the statement of
Theorem 2 there is no quasiregularity hypothesis and no restriction $p\neq2$:
the paper's exclusion of $p=2$ at paper.tex:230 comes \emph{after} the Remark
that computes the case $p=2$ explicitly, so it governs Sections 3--6 and not
the theorem.
\end{theorem}
\begin{proof}
\emph{Step 0: the constant.}  Put
\[
  K=\max\Bigl\{p-1,\ \frac1{p-1}\Bigr\} .
\]
Since $p>1$ we have $p-1>0$; if $p\ge2$ then $K\ge p-1\ge1$, and if $p<2$ then
$0<p-1<1$ so $K\ge1/(p-1)>1$.  Thus $K\ge1$, and
\noderef{EnergyIdentity} applies with this $n$, $p$ and $K$; let $C>0$ be the
constant it provides.  It depends only on $n$, $p$ and $K$, and $K$ depends only
on $p$, so $C=C(n,p)$.  Now fix $H$ satisfying (H1)--(H3).

\emph{Step 1: $H$ is $K$-quasiregular on $\mathbb{B}$.}  By
\noderef{Eq18ImpliesQuasiregular}, applied with $\Omega=\mathbb{B}$, using
$H\in W^{1,2}_{\mathrm{loc}}(\mathbb{B})$ from (H1) and (H2),
\begin{equation}\tag{1}
  H\ \text{is}\ K\text{-quasiregular on }\mathbb{B}
  \ \text{in the sense of \noderef{QuasiregularOn}} .
\end{equation}

\emph{Step 2: the coefficients.}  Let $j\in\mathbb{Z}$.  By part (a) of
\noderef{FourierCoefRegularity}, applied with the three clauses of (H1), the
function $a_j$ is continuous on $[0,1)$; and by the third clause of
\noderef{FourierCoefOde}, applied with $n\ge1$, $1<p<\infty$, (H1) and (H2),
it is $C^\infty$ on $(0,1)$.  By the second clause of \noderef{FourierCoefOde},
applied with the same data, $a_j$ satisfies
\[
  r\bigl(r\,a_j'\bigr)'(r)+pn\,r\,a_j'(r)+(n^2-j^2)(p-1)\,a_j(r)=0
  \qquad\text{for every }r\in(0,1),
\]
which is the Euler equation (35).  These are exactly the hypotheses of
\noderef{FourierCoefForm36} at the index $j$, which imposes no restriction on
$p$.  Its two cases give:
\begin{equation}\tag{2}
  |j|<n\ \Longrightarrow\ a_j(r)=0\ \text{ for every }r\in[0,1),
\end{equation}
\begin{equation}\tag{3}
  |j|\ge n\ \Longrightarrow\ \exists\,A_j\in\mathbb{C}:\
    a_j(r)=A_j\,r^{\lambda_{|j|}}\ \text{ for every }r\in[0,1) .
\end{equation}
Choose, for each $j$ with $|j|\ge n$, such a constant $A_j$, and put $A_j=0$
for $|j|<n$.  This defines $A\colon\mathbb{Z}\to\mathbb{C}$, and since
$(0,1)\subseteq[0,1)$ and $\lambda_j=\lambda_{|j|}$, (3) is exactly assertion
(b).

\emph{Step 3: proof of (39).}  The hypotheses of
\noderef{CoefCompatAndSeries} are met: $n\ge1$, $1<p<\infty$, $H$ continuous on
$\mathbb{B}$, in $W^{1,2}_{\mathrm{loc}}(\mathbb{B})$ and differentiable at
almost every point of $\mathbb{B}$, all three by (H1), $H$ a solution of (18)
by (H2), and the
coefficient hypothesis is (b) of Step 2.  Its second and
third conclusions give
\begin{equation}\tag{4}
  A_{-n}=0,\qquad A_{-k}=\epsilon_k\overline{A_k}\quad\text{for every integer }k>n .
\end{equation}
It remains to treat $k=n$.  By part (23) of \noderef{LambdaEpsEstimates},
$\epsilon_n(n,p)=0$, so $\epsilon_n\overline{A_n}=0=A_{-n}$ by the first part of
(4).  Hence
\begin{equation}\tag{5}
  A_{-k}=\epsilon_k\overline{A_k}\qquad\text{for every integer }k\ge n,
\end{equation}
which is (39).

\emph{Step 4: proof of (22).}  Apply \noderef{EnergyIdentity} with this $H$ and
this $A$.  Its hypotheses hold: (H0) there is our (H2); (H1) there is the
continuity and the almost-everywhere differentiability of our (H1); its quasiregularity hypothesis is (1);
its integrability hypothesis is our (H3); its hypothesis (H4) is assertion (b)
of Step 2; and its hypothesis (H5), namely $A_{-k}=\epsilon_k\overline{A_k}$
for $k>n$, is contained in (5).  Its clause (22a) is precisely the pair of assertions
\begin{equation}\tag{6}
  \sum_{m=0}^{\infty}(n+1+m)|A_{n+1+m}|^2\ \text{converges},\qquad
  \sum_{m=0}^{\infty}(n+1+m)|A_{n+1+m}|^2
  \le C\iint_{\mathbb{B}}\bigl(|H_\xi|^2+|H_{\bar\xi}|^2\bigr)d\sigma ,
\end{equation}
which is (22); the reindexing $k=n+1+m$ turns the left side into
$\sum_{k\ge n+1}k|A_k|^2$.

\emph{Step 5: summability of $\sum_{k\ge n}k|A_k|^2$.}  By the first part of
(6) the series $\sum_{m\ge0}(n+1+m)|A_{n+1+m}|^2$ converges.  Adding the single
further term $n|A_n|^2$,
\begin{equation}\tag{7}
  \sum_{m=0}^{\infty}(n+m)|A_{n+m}|^2
  =n|A_n|^2+\sum_{m=0}^{\infty}(n+1+m)|A_{n+1+m}|^2<\infty ,
\end{equation}
so the series of nonnegative terms $\sum_{m}(n+m)|A_{n+m}|^2$ converges, which
is hypothesis (S) of \noderef{Thm2SeriesW12Convergence}.

\emph{Step 6: the series and its sum $H'$.}  By (5) and (7) the two hypotheses
(C) and (S) of \noderef{Thm2SeriesW12Convergence} hold for this $A$.  That node
therefore provides a function $H'$ which is continuous on $\mathbb{B}$ and
$C^\infty$ on $\mathbb{B}-\{0\}$, whose Fourier coefficients satisfy
\begin{equation}\tag{8}
  a_j(H')(r)=A_j\,r^{\lambda_j}\ \ (|j|\ge n),\qquad
  a_j(H')(r)=0\ \ (|j|<n),\qquad r\in(0,1),
\end{equation}
and for which the series (19) converges to $H'$ in $W^{1,2}(\mathbb{B})$ in the
sense of its clause (g).

\emph{Step 7: $H=H'$ on $\mathbb{B}$.}  Fix $r\in(0,1)$ and consider the
$2\pi$-periodic continuous function
\[
  g_r(\theta)=H\bigl(re^{i\theta}\bigr)-H'\bigl(re^{i\theta}\bigr),
\]
continuous because both $H$ and $H'$ are continuous on $\mathbb{B}$ and
$\theta\mapsto re^{i\theta}$ is continuous with values in $\mathbb{B}$.  By
\noderef{FourierCoefOfH}, the $j$-th Fourier coefficient of $g_r$ in the
relevant normalisation is $a_j(H)(r)-a_j(H')(r)$ for every $j\in\mathbb{Z}$.
For $|j|\ge n$ this is $A_jr^{\lambda_j}-A_jr^{\lambda_j}=0$ by Step 2 (b) and
(8); for $|j|<n$ it is $0-0=0$ by (2) and (8).  So all Fourier coefficients of
the continuous $2\pi$-periodic function $g_r$ vanish, and therefore
$g_r\equiv0$, by the uniqueness theorem for Fourier coefficients of continuous
periodic functions.  As $r\in(0,1)$ was arbitrary and every $\xi$ with
$0<|\xi|<1$ is of the form $re^{i\theta}$ with $r=|\xi|\in(0,1)$, we get
$H=H'$ on $\mathbb{B}-\{0\}$.  At the origin, both $H$ and $H'$ are continuous
on $\mathbb{B}$, so
\[
  H(0)=\lim_{\xi\to0}H(\xi)=\lim_{\xi\to0}H'(\xi)=H'(0),
\]
the limits being along $\mathbb{B}-\{0\}$.  Hence $H=H'$ on $\mathbb{B}$.

\emph{Step 8: proof of (19).}  Since $H=H'$ on the open set $\mathbb{B}$, for
every $\xi\in\mathbb{B}$ the two functions agree on a neighbourhood of $\xi$,
so their real Fréchet derivatives agree there and hence, by \noderef{DZ} and
\noderef{DZbar}, $H_\xi=H'_\xi$ and $H_{\bar\xi}=H'_{\bar\xi}$ on
$\mathbb{B}-\{0\}$, where both are differentiable.  Therefore the three
integrands of clause (g) of Step 6 coincide with those of the assertion (19)
almost everywhere on $\mathbb{B}$ --- everywhere except possibly at the single
point $0$, which is a Lebesgue null set --- so the three integrals are equal
for every $N$, and the convergence to $0$ asserted in (19) is clause (g).

Steps 3, 2, 4 and 8 give (39), (b), (22) --- both its convergence assertion and
its inequality --- and (19) respectively, with the
constant $C$ of Step 0, which completes the proof.
\end{proof}
